\chapter{研究对象与问题定义}

\section{研究对象}

Nous Wave 研究跨模型、会话、进程、检索索引和外部资源持续演化的主体认知。系统需要保留可追溯的经历与认知修订，在检索和上下文组装时说明信息来源，并在更换模型或服务后端、纠正数据和恢复运行时维持权限归属一致。\cite{nous-research-package}

主体（Subject）是跨模型和运行进程延续的逻辑身份；模型（Model）提供生成能力。模型、供应商、会话和进程可以变化，主体的认知对象、证据谱系和当前状态由相应的语义所有者管理。核心区分为 \texttt{Subject $\neq$ Model}。

Nous Wave 将记忆（Memory）作为基础认知领域，并围绕主体长期连续性组织可组合能力。完整目标还包括自我、关系、动机、认知运行时、外部资源引用和跨领域离线认知。普通智能体或检索增强生成系统可组合使用记忆、运行时、检索和资源访问；需要长期主体连续性的系统可进一步启用自我、关系和动机。

\section{对话历史、检索与长期主体状态}

对话记录保存系统接收过的消息；记忆语义则说明主体当前接受什么认知、认知源自何次经历以及后续如何修订。检索增强生成可从外部知识中找到相关片段，主体是否接触材料、形成观察事件并接受长期认知由相应流程记录。向量相似度用于筛选候选材料；主体身份、认知指向、冲突处理、当前可用性和强化规则由领域合同决定。

提示词中的人格设定用于描述一次模型调用条件。长期主体状态还需分别保存自我修订、证据谱系、近期活动、任务上下文和可重建检索索引。外部知识库继续承担全文解析、索引和同步；Nous Wave 维护资源身份、访问范围和引用，并在主体实际接触材料后记录观察事件（ObservationOccurrence）。\cite{nous-research-package,mem0-commit,graphiti-commit}

对话历史和检索增强生成适用于消息回顾与外部知识查找。涉及主体特定状态、认知修订、来源核验和跨运行恢复时，系统还需相应的语义所有者、持久化记录和校验流程。

\section{研究问题}

当前研究聚焦以下问题：

\begin{enumerate}
\item 如何在主体身份与模型、会话分离时，表达来源、观察事件、证据与认知修订？
\item 如何将记忆、经历片段、日记、自我、关系、动机和运行时状态分配给语义所有者？
\item 如何融合精确路径、实体、词项、稠密向量、时间、运行时、拓扑和外部资源证据，同时保持认知权威与检索服务各自职责？
\item 如何在离线维护中提出、复核并路由假设，由领域所有者处理认知提案？
\item 如何在权威数据、运行时状态和服务投影分离后，实现并发保护、投影重建、撤销和清除？
\end{enumerate}

\section{研究状态标记}

本文以 2026-09-22 研究资料包作为当前目标语义基线，以代码快照记录实现状态。\StatusTag{目标}{由研究资料包支持的目标语义}；\StatusTag{已实现}{源树中可定位的实现}；\StatusTag{已核验}{注明平台、日期和范围的运行记录}。候选算法、局部实验观察和开放问题分别标注。

研究将长期主体连续性作为软件系统问题。心理学与哲学文献用于厘清问题结构，并与具体的数据结构和软件领域划分保持区分。\cite{olson2023,mcadams2006}
